// Halftone — rotated dot-pattern halftone effect
//
// Converts the image to a dot-pattern halftone by overlaying a rotated
// sinusoidal grid. Darker areas produce larger dots (more black), lighter
// areas produce smaller dots (more white). The grid is rotated by the
// angle parameter using sin/cos coordinate transform.
//
// Connect @image to any RGB input.

f$dot_size = 8.0;       // dot pitch in pixels (smaller = finer pattern)
f$angle = 45.0;         // grid rotation in degrees (45 = classic halftone)
f$hardness = 2.0;       // dot edge sharpness (1 = soft, 5 = crisp)
f$ink = 1.0;            // ink darkness (0 = faint, 1 = full black)

vec3 col = @image;
float lum = luma(col);

// ── Rotate UV coordinates for angled grid ───────────────────────────
float rad = $angle * PI / 180.0;
float ca = cos(rad);
float sa = sin(rad);

// Work in pixel space for consistent dot sizing
float px_x = u * iw;
float px_y = v * ih;

// Rotate around image center
float cx = iw * 0.5;
float cy = ih * 0.5;
float rx = (px_x - cx) * ca - (px_y - cy) * sa;
float ry = (px_x - cx) * sa + (px_y - cy) * ca;

// ── Generate dot pattern ────────────────────────────────────────────
// Two perpendicular sine waves create a 2D dot grid.
// Their product forms circular-ish dots at intersections.
float freq = TAU / $dot_size;
float gx = sin(rx * freq) * 0.5 + 0.5;
float gy = sin(ry * freq) * 0.5 + 0.5;

// Combine into radial dot shape (multiply creates round dots)
float dot = gx * gy;

// ── Threshold against luminance ─────────────────────────────────────
// Dark pixels need more ink = lower threshold for black
// The smoothstep controls dot edge softness
float threshold = lum;
float edge = 0.5 / $hardness;
float halftone = smoothstep(threshold - edge, threshold + edge, dot);

// ── Output ──────────────────────────────────────────────────────────
// Blend between ink color (black) and paper (white)
vec3 ink_color = vec3(1.0 - $ink);
@OUT = lerp(ink_color, vec3(1.0), halftone);
